\section{The main theorem}
\label{sec:main}

We now derive the combinatorial statement of Gilmer's theorem from the
information-theoretic strengthening of \Cref{sec:theorem1}.

\begin{theorem}[Gilmer, {\cite[Theorem 2]{gilmer2022}}]\label{thm:main}
There exists a universal constant $c > 0$ such that for every union-closed
family $\Fcal \subseteq 2^{[n]}$ with $|\Fcal| \ge 2$, some element
$i \in [n]$ satisfies
\begin{align}\label{eq:main}
\bigl| \{ A \in \Fcal : i \in A \} \bigr| \;\ge\; c \cdot |\Fcal|.
\end{align}
The explicit value $c = 0.01$ is admissible.
\end{theorem}

\begin{proof}
The strategy is contrapositive: we show that if every element $i \in [n]$
were contained in strictly fewer than a $0.01$ fraction of the sets in
$\Fcal$, then the information-theoretic strengthening
(\Cref{thm:entropy-gap}) applied to two iid uniform samples from $\Fcal$
would yield $\H(A \cup B) > \H(A)$ when $\H(A) > 0$ -- contradicting the
fact that $A \cup B \in \Fcal$ a.s.\ (so $\H(A \cup B) \le \log |\Fcal| = \H(A)$).

\paragraph{Step 1: Set up the uniform coupling on $\Fcal$.}
Let $A$ and $B$ be independent samples from the uniform distribution on
$\Fcal$. Then $A$ and $B$ each have entropy
\begin{align}\label{eq:HA-uniform}
\H(A) \;=\; \H(B) \;=\; \log |\Fcal|,
\end{align}
by definition of the uniform distribution on a set of size $|\Fcal|$.

\paragraph{Step 2: Union-closedness gives $\H(A \cup B) \le \log|\Fcal|$.}
Because $\Fcal$ is union-closed (\Cref{def:union-closed}), $A \cup B \in \Fcal$
almost surely. Hence $A \cup B$ is a random variable supported on the finite
set $\Fcal$, and \Cref{fac:uniform-max} yields
\begin{align}\label{eq:cup-bound}
\H(A \cup B) \;\le\; \log |\Fcal| \;=\; \H(A),
\end{align}
where the equality uses Eq.~\eqref{eq:HA-uniform}.

\paragraph{Step 3: Contrapositive setup.}
We argue by contradiction. Suppose, toward a contradiction, that every
$i \in [n]$ satisfies
\begin{align}\label{eq:contra-hyp}
\bigl| \{ A \in \Fcal : i \in A \} \bigr| \;<\; 0.01 \cdot |\Fcal|.
\end{align}
Since $A$ is uniform on $\Fcal$, $\Pr[i \in A] = | \{ A \in \Fcal : i \in A \} | / |\Fcal|$.
Therefore Eq.~\eqref{eq:contra-hyp} translates to
\begin{align}\label{eq:contra-prob}
\Pr[i \in A] \;<\; 0.01 \qquad \text{for every } i \in [n].
\end{align}
In particular, $\Pr[i \in A] \le 0.01$ holds for every $i$, and the
hypotheses of \Cref{thm:entropy-gap} are satisfied.

\paragraph{Step 4: $\H(A) > 0$ from $|\Fcal| \ge 2$.}
Since $|\Fcal| \ge 2$ and $A$ is uniform on $\Fcal$,
\begin{align}\label{eq:HA-positive}
\H(A) \;=\; \log |\Fcal| \;\ge\; \log 2 \;=\; 1 \;>\; 0.
\end{align}

\paragraph{Step 5: Derive the contradiction.}
Applying \Cref{thm:entropy-gap},
\begin{align}\label{eq:apply-thm-entropy}
\H(A \cup B) \;\ge\; 1.26 \, \H(A).
\end{align}
Combining Eq.~\eqref{eq:apply-thm-entropy} with Eq.~\eqref{eq:cup-bound},
\begin{align*}
1.26 \, \H(A) \;\le\; \H(A \cup B) \;\le\; \H(A),
\end{align*}
which forces $0.26 \, \H(A) \le 0$, i.e.\ $\H(A) \le 0$. This contradicts
Eq.~\eqref{eq:HA-positive}.

\paragraph{Conclusion.}
The supposition Eq.~\eqref{eq:contra-hyp} fails for at least one $i \in [n]$;
that is, some $i \in [n]$ satisfies
$|\{A \in \Fcal : i \in A\}| \ge 0.01 \cdot |\Fcal|$. Taking $c = 0.01$
gives Eq.~\eqref{eq:main}. This completes the proof.
\end{proof}

\begin{remark}\label{rem:tightness}
The constant $c = 0.01$ in \Cref{thm:main} is not tight. Gilmer's original
paper~\cite{gilmer2022} already notes that the constant was not optimized
and that the same proof technique, with a more careful analysis of the
constants in \Cref{lem:gap}, would yield a strictly larger admissible $c$.
Subsequent work (Sawin; Chase-Lovett; Alweiss-Huang-Sellke; all in late 2022,
within days of Gilmer's preprint) showed that the entropy approach can be
pushed to $c \to (3 - \sqrt{5})/2 \approx 0.38$, which is the maximum
achievable by this technique without further input. The Frankl
conjecture~\cite{frankl1995extremal}, which asks whether $c = 1/2$ is
admissible, remains open.
\end{remark}
